\part{Mininet, SDN, Open vSwitch and networking foundations}

\section{Why Mininet and Open vSwitch?}

\begin{qa}{206}{Why Mininet instead of NS-3, OMNeT++, a hardware testbed or a physical 5G network?}
\why{The claims concern real queueing disciplines (HTB, fq\_codel), real TCP congestion control and real control
commands. These are kernel behaviours a discrete-event simulator re-implements.}
\choice{Mininet: real Linux network stack, real \code{tc}, real applications, cheap to repeat 340 times, and the same
commands run on real OVS.}
\alt{ns-3 (exact virtual time, scales, but models TCP and AQM rather than running them); a hardware testbed (realistic,
costly, slow to repeat); a 5G testbed (radio realism; not available).}
\trade{Real software behaviour and fast iteration, at the cost of timing fidelity under CPU contention and no radio.}
\ev{340 runs; real-OVS replication; propagation floors exact to $<$0.01\,ms.}{\Lthree}
\end{qa}

\begin{qa}{207}{What does Mininet emulate, and what does it not?}
\textbf{Emulates:}
\begin{itemize}
\item hosts as processes in network namespaces;
\item links as veth pairs with tc shaping and delay;
\item switches (OVS or Linux bridge).
\end{itemize}
\textbf{Does not emulate:}
\begin{itemize}
\item CPU isolation (all nodes share the host CPUs);
\item NIC hardware and line-rate timing;
\item radio;
\item wide-area loss processes beyond netem.
\end{itemize}
\end{qa}

\begin{qa}{208}{How does Mininet use Linux network namespaces?}
Each host is a shell started with \code{mnexec -n} in its own network namespace: its own interfaces, routes, ARP table
and sockets, but a shared kernel, filesystem and clock. Switches stay in the root namespace. The sudo-free harness
wraps the whole emulation in one user+net+mount namespace.
\end{qa}

\begin{qa}{209}{What is the role of virtual Ethernet pairs?}
A veth pair is a virtual cable: frames written to one end appear at the other. One end sits in the host's namespace,
the other is attached to a switch port. Qdiscs on either end shape that direction.
\end{qa}

\begin{qa}{210}{How are hosts, switches, links and controllers represented?}
\begin{itemize}
\item Hosts: \code{Host} (namespaced shells).
\item Switches: \code{OVSBridge} (OVS bridge in standalone fail-mode) or \code{LinuxBridge}.
\item Links: \code{TCLink} (veth + tc).
\item Controller: none, because OVSBridge forwards with the NORMAL action (MAC learning) and needs no OpenFlow
controller. This fixed a crash, since the reference controller binary is not installed.
\end{itemize}
\end{qa}

\begin{qa}{211}{What is Open vSwitch, and why is it appropriate?}
A production-grade programmable software switch: a kernel datapath with a userspace slow path, configured through
OVSDB, programmed by OpenFlow, with QoS queues (\code{linux-htb}). It is the data plane of many SDN and NFV
deployments, so OVS-validated behaviour transfers to real virtual switching.
\end{qa}

\begin{qa}{212}{What is the difference between a Linux bridge and Open vSwitch?}
A Linux bridge is a simple L2 learning switch with no flow tables. OVS has OpenFlow match-action tables, a flow cache,
OVSDB configuration and QoS integration. Relevant behavioural differences we \emph{measured}:
\begin{itemize}
\item carrier-loss propagation to the peer is slower on OVS (recovery 1.04\,s before the event-driven fix, about
0.09\,s after);
\item MAC relearning makes the OVS outage 0.3--0.5\,s vs 0.05\,s on Linux bridges.
\end{itemize}
\end{qa}

\begin{qa}{213}{What is the difference between the SDN control plane and the data plane?}
The control plane decides how traffic should be treated (rules, routes, allocations); the data plane forwards and
queues packets. In CARS-SF the controller is our Python loop and the data plane is the kernel or OVS datapath with
its qdiscs.
\end{qa}

\begin{qa}{214}{How does an OpenFlow rule affect packet forwarding?}
A rule = match fields + priority + actions. A packet is matched against the table, the highest-priority match wins,
and its actions run (output, set\_queue, modify, drop). In the OVS QoS backend:
\code{priority=20,ip,nw_src=<UE>,actions=set_queue:q,normal} places the packet in queue $q$ and then forwards it
normally.
\end{qa}

\begin{qa}{215}{How are flow-table entries matched and prioritised?}
By priority, then (with equal priority) undefined order, so overlaps should be avoided. In our design, app rules
(priority 30) beat host rules (priority 20), and the default NORMAL flow catches the rest. The tc u32 filters follow
the same idea: flow (prio 1) $>$ app (10) $>$ host (20).
\end{qa}

\begin{qa}{216}{How are packets handled when no rule matches?}
In OpenFlow, the table-miss entry decides (send to the controller, or drop). In our standalone OVSBridge the default
rule is NORMAL. In tc, unmatched packets go to the HTB \code{default} class: general (1:40), or eMBB in 3GPP mode.
\end{qa}

\begin{qa}{217}{How do you verify that the intended forwarding rule is actually installed?}
\code{ovs-ofctl dump-flows} (collected in telemetry on OVS) and \code{tc -s class show} / \code{tc filter show}.
Unit tests assert the exact commands (\code{test_qos.py}).
\honest{The controller does not read back and compare after each change (Q139).}
\end{qa}

\begin{qa}{218}{How do you distinguish a control-plane decision from actual data-plane behaviour?}
By measuring the data plane independently: endpoint captures give per-packet delivery and delay, and HTB counters give
per-class bytes and drops. A decision counts only if those outcomes change (e.g.\ the ablation shows ceiling changes
moving p95 from \NAblSchoolFullHc{} to \NAblSchoolFullStrictHc\,ms).
\end{qa}

\begin{qa}{219}{What are the limitations of a virtual testbed on one machine?}
Shared CPUs (contention distorts timers), no hardware queues, a limited aggregate rate, all clocks identical (good for
one-way delay, unrealistic for synchronisation issues), and no radio.
\end{qa}

\begin{qa}{220}{How can host CPU scheduling distort latency and throughput measurements?}
Starved processes send late (burstier); starved kernel timers release netem and HTB packets late (inflated delay);
softirq backlog adds jitter.
\ev{Observed: a busy-waiting generator pushed load to 40--56 on 20 cores. All 111 affected runs were archived; the
generator now sleeps (\NGenCpuBefore\%$\rightarrow$\NGenCpuAfter\% CPU per 40\,Mb/s flow), with 4 parallel workers
keeping load below the core count.}{\Lthree}
\end{qa}

\begin{qa}{221}{How do virtual link parameters translate into tc configuration?}
\code{TCLink(bw, delay)} installs HTB (rate) with a netem child (delay) at the root of each veth end. Our QoS enforcer
\emph{replaces} the root qdisc on both bottleneck ends, which would silently discard any netem delay there. That is
why propagation delay is emulated on the destination-facing ports instead (§\ref{sec:delayplace}).
\end{qa}

\begin{qa}{222}{Which aspects of radio propagation, scheduling, handover and interference are absent?}
All of them: no PHY or MAC, no PRB scheduling, no HARQ, no fading, no mobility. The 5G deployment adds the
\emph{core} (NAS, PDU sessions, slices) but UERANSIM simulates the RAN without radio impairments.
\end{qa}

\begin{qa}{223}{What conclusions can you legitimately draw from Mininet experiments?}
Relative comparisons of allocation, AQM and healing policies at a shared IP bottleneck with real kernel queueing and
real TCP; the correctness and timing of the control loop; measurement-method validity.
\textbf{Not:} absolute field latencies, radio-limited behaviour, scalability to large networks.
\end{qa}

\begin{qa}{224}{Which conclusions need a radio-aware simulator or a real 5G testbed?}
Radio scheduling interaction (PRB allocation per slice), coverage and handover effects, wireless loss, RAN-slice
isolation (as in the isolation paper [2]), and end-to-end latency budgets including the air interface.
\end{qa}

\section{Topology and resource modelling}
\label{sec:delayplace}

\FigTopo

\begin{qa}{225}{Why this rural topology?}
It is the smallest topology containing all three phenomena under study:
\begin{itemize}
\item one shared rural uplink (contention);
\item two service tiers (edge MEC vs core cloud), so that failures affect services unequally;
\item a primary and a weaker standby backhaul (healing).
\end{itemize}
\end{qa}

\begin{qa}{226}{What do the switches, hosts, links and backup paths represent?}
\begin{itemize}
\item \code{s_rural}: village aggregation (cell site / gNB aggregation).
\item \code{s_edge}: edge cloud (MEC, clinic and agri servers).
\item \code{s_core}: core cloud (education and internet services).
\item UEs: 2 per service institution.
\item The primary backhaul is fiber; the backup is microwave or satellite-class (50\,Mb/s, 40\,ms).
\end{itemize}
\end{qa}

\begin{qa}{227}{Which links are bottlenecks?}
Only the rural uplink (100\,Mb/s, both directions). Access links are deliberately unshaped. On backup, the 50\,Mb/s
backup link bounds core-bound traffic, and CARS-SF caps core slices to 95\% of it.
\end{qa}

\begin{qa}{228}{How are link capacities, delays, loss rates and queue limits configured?}
\begin{itemize}
\item Capacity: HTB root 100\,Mb/s (swept 50--150).
\item Delay: netem on destination-facing ports (8\,ms edge path, 28\,ms core path, backup +40\,ms).
\item Loss: zero by default; injected by the degradation scenario (+0.3\% per tick).
\item Queues: fq\_codel defaults (target 5\,ms, interval 100\,ms, limit 10240 packets); the no-slicing baseline uses
pfifo 1000.
\end{itemize}
\end{qa}

\begin{qa}{229}{Are the topology parameters measured, from the literature, or hypothetical?}
Hypothetical but plausible (fiber vs microwave delays, a 100\,Mb/s rural uplink), and swept for sensitivity.
\honest{Not calibrated against a specific deployment.}
\end{qa}

\begin{qa}{230}{How do you verify the topology after initialisation?}
A connectivity ping (healthcare UE $\rightarrow$ edge) on start; the first-tick RTT probes for every service; and,
post-run, the measured propagation floors equal configuration exactly (\NFloorEdge{} and \NFloorCore\,ms).
\end{qa}

\begin{qa}{231}{How do you identify a failed link and distinguish it from a slow link?}
\begin{itemize}
\item Failed: carrier loss (operstate) on \code{s_edge-bh}, seen by the 50\,ms link watcher.
\item Slow: rising RTT trend for core-bound probes with flat edge probes and \emph{no} local congestion evidence
(predictor).
\item Congested: drops, backlog or saturation at the local bottleneck.
\end{itemize}
\end{qa}

\begin{qa}{232}{What happens if a switch fails instead of a link?}
Core-switch failure scenario: all \code{s_core} ports go down. Carrier loss is seen as a cut, so the engine switches
to backup; the backup also ends on \code{s_core}, so the core stays unreachable. After 2 ticks the diagnosis is
``core switch failure'' (\textsf{DEGRADED}): core slices are held at their floors and edge slices get the capacity.
Diagnosed in \NCoreDiagnosedShare{} of runs.
\end{qa}

\begin{qa}{233}{What if the primary and backup routes share a physical failure domain?}
Then the backup does not help, which is exactly our core-switch case (both paths terminate on \code{s_core}). The
system localises the fault and degrades gracefully, but cannot restore core services.
Lesson: logical redundancy $\ne$ physical diversity.
\end{qa}

\begin{qa}{234}{How do you prevent an alternate route from creating a Layer-2 loop?}
By construction, with no STP:
\begin{enumerate}
\item the backup is admin-down on the \code{s_edge} side while the primary is up;
\item on failure, the primary is \emph{quarantined} (admin-down) \emph{before} the backup comes up;
\item on fail-back, the backup goes down \emph{before} the primary comes up;
\item only \code{s_edge}-side ports are driven;
\item the healer ignores link events caused by its own actions.
\end{enumerate}
\ev{Unit test asserting the exact ordering; no broadcast storms in 27 resilience runs plus the real-OVS runs.}{\Lthree}
\end{qa}

\begin{qa}{235}{How do MAC learning, ARP and forwarding state affect failover?}
After a path change, switches still map MACs to the dead port until they relearn. ARP caches are unaffected (same
IP/MAC pairs). Mitigation: flush the FDB on all three switches (\code{ovs-appctl fdb/flush},
\code{bridge fdb flush}). The residual OVS outage (0.3--0.5\,s) is mostly relearning.
\end{qa}

\begin{qa}{236}{How do you handle stale forwarding entries after a path change?}
An explicit FDB flush on activation and fail-back. Linux bridges also flush entries on a port that loses carrier.
\end{qa}

\begin{qa}{237}{How do you prove the backup route works before switching traffic to it?}
\honest{We do not pre-verify the backup. Switching happens, then reachability is checked: if the core is still
unreachable after 2 ticks, the engine escalates to \textsf{DEGRADED}.}
Better: continuous OAM on the standby (periodic probes over an isolated VLAN, or BFD on the standby link), so a dead
backup is known before it is needed.
\end{qa}

\begin{qa}{238}{Logical redundancy vs physically independent redundancy?}
Logical: two paths in the topology graph. Physical: independent conduits, power and end equipment. Our backup is
logically redundant but shares the \code{s_core} failure domain (Q233).
\end{qa}

\begin{qa}{239}{How would you model a power failure affecting multiple devices?}
Take a set of nodes down at once (all ports of the affected switches), as in our core-switch injection, possibly with
correlated recovery times. Not implemented as a scenario.
\end{qa}

\begin{qa}{240}{How would you model partial connectivity rather than complete failure?}
Backhaul degradation: netem delay +15\,ms and loss +0.3\% per tick, then a cut. \code{inject_latency} also exists.
Further options: rate reduction, intermittent flapping, asymmetric loss.
\end{qa}

\begin{qa}{241}{How do you prevent a topology change from invalidating previously collected telemetry?}
Counters are read per interface name and converted to deltas with negative values clamped to 0, so a reset or
recreated interface yields one zero sample, not garbage. Datasets carry \code{run_id} and are never mixed across
topologies without that key.
\end{qa}

\begin{qa}{242}{How would the topology scale to more villages, gateways, services and bottlenecks?}
The builder is parametric in UEs per service (tested at 1, 2 and 4) and in capacity. Multiple villages need multiple
\code{s_rural} instances, per-uplink QoS and a regional controller (Q130).
\honest{Single-machine emulation will run out of CPU at around tens of nodes with this traffic.}
\end{qa}

\section{TCP/IP, transport behaviour and traffic control}

\begin{qa}{243}{Why is TCP behaviour relevant to your results?}
Elastic TCP flows (library, web, DASH) adapt to whatever share they get, so their demand is not directly observable,
and they fill queues (bufferbloat) unless AQM intervenes. The demand estimator's elastic probe-up and the fq\_codel
choice both exist because of TCP.
\end{qa}

\begin{qa}{244}{How would TCP congestion control respond to loss or increased delay?}
Linux default CUBIC: multiplicative decrease of cwnd on loss, cubic regrowth, largely insensitive to delay until
loss. fq\_codel turns standing queues into early drops (or ECN marks), so TCP backs off before delay builds.
\end{qa}

\begin{qa}{245}{Why might throughput fall even when nominal capacity is unchanged?}
\begin{itemize}
\item rising RTT (backhaul degradation) slows cwnd growth;
\item random loss (degradation) triggers backoff;
\item receiver window limits;
\item HTB ceilings;
\item competition from unresponsive UDP in a FIFO queue.
\end{itemize}
\end{qa}

\begin{qa}{246}{How does UDP differ from TCP under congestion?}
UDP does not back off. In a FIFO queue it can starve TCP and inflate delay for everyone. HTB per-slice classes and
fq\_codel per-flow queues isolate it.
\end{qa}

\begin{qa}{247}{How can retransmissions distort application-level latency measurements?}
A retransmitted segment's delivery time includes timeout or fast-retransmit delay that a per-packet one-way delay
does not show. We match each transmission separately (flagging retransmissions), so our one-way delay is network
delay, not application latency.
\honest{Application-level latency for TCP (request to response) is not measured.}
\end{qa}

\begin{qa}{248}{What is the difference between RTT and one-way latency?}
RTT is the sum of both directions plus host turnaround. One-way delay is a single direction and needs synchronised
clocks. Our endpoints share one kernel clock, so the packet-matched one-way delay is exact (validated by the floors).
Online probes measure RTT, and monitoring uses RTT/2.
\end{qa}

\begin{qa}{249}{What is jitter, and how does its definition affect the result?}
We report the mean absolute difference of consecutive one-way delays per window (RFC 3550-style, without the
1/16 smoothing). Other definitions (standard deviation of delay, p95$-$p50) give different magnitudes; we report one
definition consistently.
\end{qa}

\begin{qa}{250}{Packet loss vs queue drops vs application-level delivery failure?}
\begin{itemize}
\item Packet loss: a packet transmitted but never received (endpoint matching).
\item Queue drops: HTB or fq\_codel drop counters (a subset of loss).
\item Application failure: data not delivered in time. For TCP, loss becomes delay through retransmission; for
real-time UDP, loss means missing data.
\end{itemize}
\end{qa}

\begin{qa}{251}{What is a Linux qdisc?}
A queueing discipline attached to an interface's egress that decides how packets are queued and dequeued.
Classless qdiscs (pfifo, fq\_codel, netem) and classful ones (HTB) form trees with filters that pick the class.
\end{qa}

\begin{qa}{252}{How does HTB enforce rate ceilings and borrowing?}
Each class has a guaranteed \code{rate} and a maximum \code{ceil}. A class under its rate sends freely; between rate
and ceil it \emph{borrows} its parent's spare tokens, in \code{prio} order and quantum-weighted; above ceil it waits.
We use rate = allocation and ceil = link capacity (work-conserving) or ceil = rate (strict).
\end{qa}

\begin{qa}{253}{How do queue limits affect delay and drops?}
A larger limit means fewer drops but more standing delay (bufferbloat): 1000 packets at 100\,Mb/s is about 120\,ms.
fq\_codel bounds sojourn time instead of length (target 5\,ms).
\ev{Without AQM, the school no-AQM ablation, the OVS queues and the 5G no-slicing mode all show large p95; the 5G
no-slicing healthcare p95 was \NFgDisasterNoneHcP\,ms.}{\Lthree}
\end{qa}

\begin{qa}{254}{What is the purpose of netem?}
It emulates delay, jitter, loss, duplication and reordering on an interface. We use it for propagation delay and
backhaul degradation.
\end{qa}

\begin{qa}{255}{What are the limits of using traffic shaping to represent network conditions?}
Shaping is deterministic and memoryless unless configured otherwise. It does not reproduce wireless loss bursts or
link-layer retransmission. Timer granularity and CPU contention add artefacts. Delay placement interacts with other
qdiscs (Q221).
\end{qa}

\begin{qa}{256}{How does FQ-CoDel differ from a simple FIFO queue?}
FQ: per-flow hashing into sub-queues with deficit round robin, so sparse flows (ECG, voice) are served ahead of bulk
flows. CoDel: drops when sojourn time stays above 5\,ms for an interval, which keeps standing queues short.
FIFO: one queue, tail drop, no isolation.
\end{qa}

\begin{qa}{257}{Why can strict priority scheduling cause starvation?}
A higher class with continuous backlog always wins, so lower classes never transmit. CARS-SF avoids pure strict
priority: each class has a guaranteed rate (floor or allocation), and \code{prio} only orders borrowing of spare
capacity.
\end{qa}

\begin{qa}{258}{How does weighted fair queuing compare with strict priority?}
WFQ shares capacity in proportion to weights (no starvation, bounded delay per weight). Strict priority minimises
high-class delay at the risk of starvation. HTB with rates is closer to WFQ with guarantees; HTB prio adds mild
priority for spare capacity.
\end{qa}

\begin{qa}{259}{What is the difference between rate limiting, shaping and policing?}
Shaping delays excess packets (queue, smooth output). Policing drops or marks excess immediately (no queue). Rate
limiting is the generic term. HTB shapes; we do not police.
\end{qa}

\begin{qa}{260}{How do you determine whether the QoS policy actually influences the intended class?}
Per-class HTB counters (bytes, drops, backlog) per direction each tick; packet-matched per-slice delay and loss;
ablations that change one QoS parameter and observe the effect; and the classification tests (u32 filter strings).
\end{qa}

\begin{qa}{261}{How do you distinguish a configured ceiling from the throughput a flow actually receives?}
The configured value is in \code{controller_history}; delivered throughput per window comes from captures. They
differ by demand (a flow may use less), by borrowing (it may use more, up to ceil), and by TCP dynamics.
\end{qa}

\begin{qa}{262}{How do you inspect and debug interfaces, routes, sockets, qdiscs and OVS rules?}
\code{ip -s link}, \code{ip route}, \code{ss -tinp}, \code{tc -s qdisc/class/filter show dev X},
\code{ovs-vsctl show}, \code{ovs-ofctl dump-flows/dump-ports}, \code{ovs-appctl fdb/show} and
\code{ofproto/trace}, \code{tcpdump -i X}, \code{bridge fdb show}. Inside Mininet: \code{mnexec -a <pid>}.
\end{qa}

\begin{qa}{263}{The controller claims a route changed, but packets still use the old path. How do you debug it?}
\begin{enumerate}
\item Check the port states (\code{ip link}) on both ends.
\item Check the counters on both links (\code{ip -s link}) to see which one carries the bytes.
\item Inspect the FDB (\code{ovs-appctl fdb/show s_edge}): stale MAC entries pointing to the old port are the usual cause.
\item Use \code{ofproto/trace} for one packet.
\item Run tcpdump on both candidate links.
\end{enumerate}
Fix: flush the FDB, or check that the quarantine actually took effect.
\end{qa}
